% ch7_b 第7章 §7.1尾–§7.2 (书页 288–300 / p303–p315)
%JOIN%
理解为什么它具有填充因子 $\nu=1/3$。这是因为 $z_i$（固定 $i$）的最高次幂是 $3(N-1)$。角动量为 $3(N-1)$ 的轨道的半径为 $r_{\max}=\sqrt{2\times 3(N-1)}\,l_B$。这里 $r_{\max}$ 也是液滴的半径，是 $\nu=1$ 液滴半径 $\sqrt{2\times(N-1)}\,l_B$ 的 $\sqrt{3}$ 倍。因此，对 $\Psi_3$ 有 $\nu=\frac{1}{3}$。

为了理解为什么 $\Psi_3$ 具有均匀的密度，让我们考虑下面电子位置的联合概率分布：
\begin{equation*}
P(z_1\cdots z_N)\propto|\Psi_m(z_1\cdots z_N)|^2=\mathrm{e}^{-\beta V(z_1\cdots z_N)}
\tag{7.2.4}
\end{equation*}
其中我们已把 $\Psi_3$ 态推广为 $\Psi_m$ 态
\begin{equation*}
\Psi_m=\prod_{i<j}(z_i-z_j)^m\,\mathrm{e}^{-\sum|z_i|^2/4l_B^2}
\end{equation*}
我们看到，为使 $\Psi_m$ 全反对称，$m$ 必须是奇整数。在式 (7.2.4) 中选择 $T=\frac{1}{\beta}=\frac{m}{2}$，我们可以把
\begin{equation*}
V=-m^2\sum_{i<j}\ln|z_i-z_j|+\frac{m}{4l_B^2}\sum_i|z_i|^2
\end{equation*}
视为由带"电荷" $m$ 的粒子组成的二维等离子体的势，而把 $P(z_1\cdots z_N)$ 视为该等离子体的概率分布。"电荷"为 $m_1$ 与 $m_2$ 的两个粒子之间的势是 $-m_1m_2\ln(r)$，力是 $m_1m_2/r$。因此，要理解 $\Psi_m$ 的密度分布，我们只需计算等离子体的"电荷"分布。

在 $V$ 中我们看到，等离子体中的每个粒子感受到一个势 $m|z|^2/4l_B^2$，它可以视为由一个背景"电荷"分布产生的势。实际上，这样的势由"电荷"密度为 $\rho_\phi=1/2\pi l_B^2$ 的均匀背景"电荷"产生，这正好就是磁通量子的密度。为看出这一点，我们注意到，这样的背景"电荷"作用在半径 $r$ 处的"电荷" $m$ 上的力，是"电荷" $m$ 与半径 $r$ 之内总"电荷"之间的力。该力由 $F=m\rho_\phi\pi r^2/r=m\rho_\phi\pi r$ 给出。这给出一个势 $\frac{m}{2}\pi\rho_\phi r^2$，正好就是上面的势。由于等离子体的完全屏蔽性质，等离子体倾向于"电荷"中性。等离子体的"电荷"密度等于背景"电荷"密度 $mn_e=\rho_\phi$，其中 $n_e$ 是电子密度。我们得到填充因子 $\nu=n_e/\rho_\phi=\frac{1}{m}$。

为什么 $\Psi_m$ 是不可压缩态？注意到 $\Psi_m$ 的总角动量（$z_i$ 的总幂次）为 $M_0=m\frac{N(N-1)}{2}$。$\Psi_m$ 是具有最小总角动量的态，它在每一对电子之间都有 $m$ 阶零点。压缩液滴会减小总角动量并产生更低阶的零点，这要耗费能量。我们预期该能隙为库仑能 $e^2/l_B$ 的量级。上述论断没有数学证明，但我们有大量的数值证据。

我们想指出，虽然对具有库仑相互作用的电子而言 $\Psi_3$ 只是近似基态，但存在一个理想哈密顿量，使 $\Psi_3$ 成为它的精确基态。理想哈密顿量具有形式
\begin{equation*}
H=\frac{1}{2m}\sum_i(\partial_{\pmb{x}_i}-\mathrm{i}q_e\pmb{A})^2-\sum_{i<j}UV(\pmb{r}_i-\pmb{r}_j),\quad V(\pmb{r})=\partial_{\pmb{r}}^2\delta(\pmb{r}).
\tag{7.2.5}
\end{equation*}
可以证明，对任意态 $\psi$ 有 $\langle\psi|H|\psi\rangle\geqslant\frac{1}{2}N\hbar\omega_c$，而对 $\Psi_3$ 有 $\langle\Psi_3|H|\Psi_3\rangle=\frac{1}{2}N\hbar\omega_c$。因此 $\Psi_3$ 确实是 $H$ 的精确基态。数值计算还表明，对理想哈密顿量而言 $\Psi_3$ 是不可压缩的。

\subsection*{问题 7.2.3}

证明：对理想哈密顿量 (7.2.5) 有 $\langle\Psi_3|H|\Psi_3\rangle=\frac{1}{2}N\hbar\omega_c$。为 $\nu=1/m$ 的 Laughlin 态 $\Psi_m$ 找一个理想哈密顿量。

\subsection*{问题 7.2.4}

双层 FQH 态由波函数
\begin{equation*}
\Psi_{lmn}=\prod_{i<j}(z_i-z_j)^l\prod_{i<j}(w_i-w_j)^m\prod_{i,j}(z_i-w_j)^n\,\mathrm{e}^{-\sum|z_i|^2/4l_B^2}\,\mathrm{e}^{-\sum|w_i|^2/4l_B^2}
\tag{7.2.6}
\end{equation*}
描述，其中 $z_i$ 是第一层中电子的坐标，$w_i$ 是第二层中电子的坐标。这样的态称为 $(lmn)$ 态。设 $N_1$ 与 $N_2$ 为两层中的电子数。求使两层中电子液滴尺寸相等的比值 $N_1/N_2$。求两层中电子的总填充因子。（提示：可先假设 $n=0$ 以简化问题。）

\subsection{具有分数电荷与分数统计的准粒子}

\begin{keypoints}
\item 由等离子体类比得到的分数电荷与分数统计。
\end{keypoints}

不可压缩基态 $\Psi_m$ 之上的准空穴激发由多体波函数
\begin{equation*}
\Psi^h(\xi,\xi^*)=\sqrt{C(\xi,\xi^*)}\prod_i(\xi-z_i)\Psi_m,
\end{equation*}
描述，其中 $\xi$ 是准空穴的位置，$C(\xi,\xi^*)$ 是归一化系数。计算准空穴电荷最简单的方式，是先从 $\Psi_m$ 态中移除一个电子，得到波函数 $\prod_i(\xi-z_i)^m\Psi_m$，它在 $\xi$ 处有一个电荷为 $e$ 的空穴。然而，$\prod_i(\xi-z_i)^m\Psi_m$ 也可以视为在 $\xi$ 处有 $m$ 个准空穴的波函数。因此，准空穴携带的电荷是 $e/m$。

%FIG{7.13}: {电子的电荷与测试电荷必须中和背景电荷。}

计算准空穴电荷更直接的途径是等离子体类比。准空穴态 $\Psi_h$ 中电子位置的分布为
\begin{equation*}
P^h_\xi(\{z_i\})=\bigl|\Psi^h(\xi,\xi^*)\bigr|^2=\mathrm{e}^{-\beta[V(z_i)-m\ln|\xi-z_i|]}
\end{equation*}
因此，电子感受到一个附加的背景势 $-m\ln|\xi-z_i|$。这个势由一个单位测试电荷产生。背景电荷在 $\xi$ 处包含一个额外的单位测试电荷。为维持电荷中性，$\xi$ 附近必须少 $1/m$ 个电子（见图 7.13）。这样，准空穴携带电荷 $e/m$。

为了理解准空穴的分数统计（Arovas et al., 1984），让我们从下面几个准粒子的有效拉格朗日量出发：
\begin{equation*}
Z=\int\mathcal{D}[\pmb{x}_i(t)]\;\mathrm{e}^{\mathrm{i}\int\mathrm{d}t\,\sum_i[-E_0+\frac{1}{2}m\dot{\pmb{x}}_i^2+\pmb{a}(\pmb{x}_1,\pmb{x}_2,\ldots)\cdot\dot{\pmb{x}}_i]}
\end{equation*}
有效规范势 $\pmb{a}$ 决定统计。正如 2.3 节讨论的自旋系统那样，$\pmb{a}$ 由相干态 $|\Psi^h(\xi(t),\xi^*(t))\rangle$ 的贝里相决定。

让我们先只考虑单个准粒子。对于准粒子 $\xi(t)$ 的绝热运动，$\Psi^h(\xi(t),\xi^*(t))$ 的相位变化给出贝里相。对小的 $\Delta t$，贝里相为
\begin{equation*}
\langle\Psi^h(\xi(t+\Delta t),\xi^*(t+\Delta t))|\Psi^h(\xi(t),\xi^*(t))\rangle=\mathrm{e}^{\mathrm{i}\Delta t\,\pmb{a}\cdot\dot{\pmb{x}}}.
\end{equation*}
于是 $\pmb{a}$ 由
\begin{equation*}
\mathrm{i}\langle\Psi^h(\xi(t),\xi^*(t))|\frac{\mathrm{d}}{\mathrm{d}t}|\Psi^h(\xi(t),\xi^*(t))\rangle=\pmb{a}\cdot\dot{\pmb{x}}
\end{equation*}
给出。由于 $\Psi^h$ 只通过 $\xi(t)$ 依赖 $t$，我们有
\begin{align*}
\pmb{a}\cdot\dot{\pmb{x}}&=\mathrm{i}\langle\Psi^h(\xi(t),\xi^*(t))|\frac{\mathrm{d}}{\mathrm{d}t}|\Psi^h(\xi(t),\xi^*(t))\rangle\\
&=\frac{\mathrm{d}\xi}{\mathrm{d}t}\,\mathrm{i}\langle\Psi^h(\xi,\xi^*)|\frac{\partial}{\partial\xi}|\Psi^h(\xi,\xi^*)\rangle+\frac{\mathrm{d}\xi^*}{\mathrm{d}t}\,\mathrm{i}\langle\Psi^h(\xi,\xi^*)|\frac{\partial}{\partial\xi^*}|\Psi^h(\xi,\xi^*)\rangle
\end{align*}
这样，如果我们写出 $\pmb{a}\cdot\mathrm{d}\pmb{x}=a_\xi\mathrm{d}\xi+a_{\xi^*}\mathrm{d}\xi^*$，就有\footnote{我们用到了
\[
\begin{aligned}
&\mathrm{i}\langle\Psi^h(\xi,\xi^*)|\frac{\partial}{\partial\xi}|\Psi^h(\xi,\xi^*)\rangle=\mathrm{i}\langle\Psi(\xi)|\frac{1}{\sqrt{C(\xi,\xi^*)}}\frac{\partial}{\partial\xi}\frac{1}{\sqrt{C(\xi,\xi^*)}}|\Psi(\xi)\rangle\\
&\quad=\mathrm{i}\int\prod_i\mathrm{d}^2z_i\;\Psi^*(\xi^*)\frac{1}{\sqrt{C}}\frac{\partial}{\partial\xi}\Bigl(\frac{1}{\sqrt{C}}\Psi(\xi)\Bigr)\\
&\quad=\mathrm{i}\frac{\partial}{\partial\xi}\int\prod_i\mathrm{d}^2z_i\;\Psi^*(\xi^*)\frac{1}{C}\Psi(\xi)-\mathrm{i}\int\prod_i\mathrm{d}^2z_i\;\Psi^*(\xi^*)\Bigl(\frac{\partial}{\partial\xi}\frac{1}{\sqrt{C}}\Bigr)\frac{1}{\sqrt{C}}\Psi(\xi)\\
&\quad=-\mathrm{i}\sqrt{C}\frac{\partial}{\partial\xi}\frac{1}{\sqrt{C}}
\end{aligned}
\]}
\begin{equation*}
a_\xi=\mathrm{i}\langle\Psi^h(\xi,\xi^*)|\frac{\partial}{\partial\xi}|\Psi^h(\xi,\xi^*)\rangle=-\mathrm{i}\sqrt{C}\frac{\partial}{\partial\xi}\frac{1}{\sqrt{C}}
\end{equation*}
这里 $C(\xi,\xi^*)=\langle\Psi(\xi)|\Psi(\xi)\rangle$，而 $\Psi(\xi)=\prod(\xi-z_i)\Psi_m$ 是 $\xi$ 的解析函数。上述结果的关键之点是，$\Psi$（$\xi$ 的解析函数）的归一化给出贝里相：
\begin{equation*}
a_\xi=+\frac{\mathrm{i}}{2}\frac{\partial}{\partial\xi}\ln C,\qquad a_{\xi^*}=-\frac{\mathrm{i}}{2}\frac{\partial}{\partial\xi^*}\ln C
\end{equation*}
归一化 $C(\xi,\xi^*)$ 可以用等离子体类比算出。我们注意到，在加入一个代表测试"电荷"与背景"电荷"之间相互作用的项 $\mathrm{e}^{-|\xi|^2/4ml_B^2}$ 之后，\footnote{注意 $\mathrm{e}^{-|\xi|^2/4l_B^2}$ 对应于一个电子与背景"电荷"之间的相互作用。一个电子对应于 $m$ 个测试"电荷"的束缚态。}
\begin{equation*}
\int\prod_i\mathrm{d}^2z_i\;\Bigl|\mathrm{e}^{-\frac{1}{4ml_B^2}|\xi|^2}\prod(\xi-z_i)\prod(z_i-z_j)^m\,\mathrm{e}^{-\frac{1}{4l_B^2}|z_i|^2}\Bigr|^2\propto\mathrm{e}^{-\beta V(\xi,\xi^*)}
\end{equation*}
给出的是在 $\xi$ 处带有一个测试"电荷"的等离子体的总能量。同样，由于等离子体的完全屏蔽性质，作用在插入的测试"电荷"上的总力为零；因而等离子体的能量不依赖于 $\xi$，即 $V(\xi,\xi^*)=\text{常数}$。我们看到
\begin{equation*}
C(\xi,\xi^*)\propto\mathrm{e}^{\frac{1}{2l_B^2}\frac{1}{m}|\xi|^2}
\end{equation*}
这给出
\begin{equation*}
a_\xi=\mathrm{i}\frac{1}{m}\frac{1}{4l_B^2}\xi^*,\qquad a_{\xi^*}=-\mathrm{i}\frac{1}{m}\frac{1}{4l_B^2}\xi
\end{equation*}
上述矢量势 $\pmb{a}$ 正比于磁场的矢量势：
\begin{equation*}
a_\xi=-\frac{1}{m}q_eA_z,\qquad a_{\xi^*}=-\frac{1}{m}q_eA_{z^*}
\end{equation*}
绕环路一周，准空穴获得 A--B 相 $\oint\mathrm{d}\pmb{x}\cdot\pmb{a}=\bigl(-\frac{q_e}{m}\bigr)\oint\mathrm{d}\pmb{x}\cdot\pmb{A}$，即一个电荷为 $-q_e/m=e/m$ 的粒子的相位。

对于两个准空穴，带两个测试"电荷"的等离子体的总能量 $V$ 由
\begin{align*}
&\mathrm{e}^{-\beta V(\xi_1,\xi_1^*,\xi_2,\xi_2^*)}\\
&\quad=\Bigl|\mathrm{e}^{-\frac{1}{4l_B^2}\frac{1}{m}(|\xi_1|^2+|\xi_2|^2)}\,|\xi_1-\xi_2|^{\frac{1}{m}}\prod_{a,i}(\xi_a-z_i)\prod(z_i-z_j)^m\,\mathrm{e}^{-\frac{1}{4l_B^2}\sum|z_i|^2}\Bigr|^2
\end{align*}
给出，其中额外的项 $|\xi_1-\xi_2|^{2/m}=\mathrm{e}^{-\beta(-)\ln|\xi_1-\xi_2|}$ 表示两个测试"电荷"之间的直接相互作用，每个"电荷"为 1。同样，$V(\xi_1,\xi_1^*,\xi_2,\xi_2^*)$ 不依赖于 $\xi_1$ 与 $\xi_2$。于是
\begin{equation*}
|\Psi(\xi_1,\xi_2)|^2\propto\mathrm{e}^{\frac{1}{2l_B^2}\frac{1}{m}(|\xi_1|^2+|\xi_2|^2)}\,|\xi_1-\xi_2|^{-\frac{2}{m}}
\end{equation*}
令 $\xi=\xi_1$、$\xi_2=0$，我们发现
\begin{align*}
a_\xi&=\mathrm{i}\frac{1}{m}\frac{1}{4l_B^2}\xi^*-\frac{\mathrm{i}}{2m}\frac{\partial}{\partial\xi}\ln|\xi|^2=\mathrm{i}\frac{1}{m}\frac{1}{4l_B^2}\xi^*-\frac{\mathrm{i}}{2m}\frac{1}{\xi}\\
a_{\xi^*}&=-\mathrm{i}\frac{1}{m}\frac{1}{4l_B^2}\xi+\frac{\mathrm{i}}{2m}\frac{1}{\xi^*}
\end{align*}
这给出 $(a_x,a_y)=(-y,x)\frac{1}{r^2}\frac{1}{m}+\cdots$。当我们把准空穴 $\xi_1$ 绕 $\xi_2$ 移动一周时，准空穴获得的相位为 $\frac{e}{m}B\times\text{面积}+\frac{2\pi}{m}$。第一项来自均匀磁场。附加项 $\frac{2\pi}{m}$ 来自 $(-y,x)\frac{1}{r^2}\frac{1}{m}$（译注：原文此处印作 $\frac{1}{r}$）。正是这一项赋予准空穴统计角为 $\theta=\frac{\pi}{m}$ 的分数统计（见式 (7.1.1)）。

\subsection*{问题 7.2.5}

考虑双层 $(lmn)$ 态，即式 (7.2.6)。第一层中的准空穴由波函数
\begin{equation*}
\Psi^h_1(\xi_1)=\prod_i(\xi_1-z_i)\Psi_{lmn}(\{z_i,w_j\})
\end{equation*}
描述，第二层中的准空穴由
\begin{equation*}
\Psi^h_2(\xi_2)=\prod_i(\xi_2-w_i)\Psi_{lmn}(\{z_i,w_j\})
\end{equation*}
描述。求第一层与第二层中准空穴的分数电荷与分数统计。求第一层准空穴与第二层准空穴之间的互统计。两个非全同粒子之间的互统计定义为把一个粒子绕另一个粒子一整周所产生的贝里相。

%FIG{7.14}: {分级 $\nu=2/7$ 与 $\nu=2/5$ 态。}

\subsection{分级分数量子霍尔态——Laughlin 理论的推广}

\begin{keypoints}
\item 分级 FQH 态等价于准粒子的 Laughlin 态。
\end{keypoints}

要构造填充因子不是 $1/m$ 的更一般的 FQH 态，我们可以从 $\nu=1/m$ 的 Laughlin 态出发，加入准空穴或准粒子来改变平均电子密度。当其密度达到某一数值时，准空穴/准粒子气体自身也能形成一个 Laughlin 态。所得的态就是一个分级 FQH 态（见图 7.14）。

当 $\nu=1/3$ 态之上的电荷为 $e/3$ 的准空穴凝聚成一个 Laughlin 态时，我们得到 $\nu=2/7$ 的 FQH 态。波函数具有形式
\begin{align*}
&\Psi(z_1\cdots z_N)\\
&\quad=\int\prod_i\mathrm{d}^2\xi_i\;\prod(z_i-z_j)^3\,\mathrm{e}^{-\frac{1}{4l_B^2}\sum|z_i|^2}\prod(\xi_i-z_j)(\xi_i^*-\xi_j^*)^2\,\mathrm{e}^{-\frac{1}{4l_B^2}\frac{1}{3}|\xi_i|^2}
\end{align*}
当 $\nu=1/3$ 态之上的电荷为 $e/3$ 的准粒子凝聚成一个 Laughlin 态时，我们得到 $\nu=2/5$ 的 FQH 态。波函数为
\begin{equation*}
\Psi_c(z_i-z_N)=\int\prod_i\mathrm{d}^2\xi_i\;\prod_{i<j}(\xi_i-\xi_j)^2(\xi_i^*-2\partial_{z_i})\prod(z_i-z_j)^3\,\mathrm{e}^{-\frac{1}{4l_B^2}\sum|z_i|^2}
\end{equation*}
（译注：原书此式末尾的指数印作 $-\frac{1}{4l_B^2}|z_i|$，疑脱求和号与平方。）

我们可以用等离子体类比计算上述两个态的性质，但计算很复杂。下一节我们将推导 FQH 态的有效场论，并用有效理论计算上述分级 FQH 态的性质。有效场论大大简化了计算，我们能容易地得到 FQH 态的填充因子与准粒子的量子数。

\section{分数量子霍尔液体的有效理论}

\begin{keypoints}
\item FQH 态的低能有效理论是 $U(1)$ Chern--Simons 理论，它们捕捉 FQH 态的普适性质。
\end{keypoints}

我们已经见过几种不同的 FQH 态。一方面，这些 QH 态具有不同的分数电荷与不同的分数统计，这暗示不同的 FQH 态属于不同的量子相。另一方面，这些 FQH 态又都具有相同的对称性，我们无法用破缺的对称性来区分它们。实际上，FQH 态包含一种新型的序——拓扑序（见第 8 章）。我们需要用新的工具来刻画拓扑序。

系统研究拓扑序的一种途径是为 FQH 液体构造低能有效理论。有效理论能够捕捉 FQH 液体的普适性质，并为如何刻画和标记 FQH 液体中不同的拓扑序提供线索。下面我们将介绍一种构造有效理论的方法，它与 Haldane 和 Halperin 提出的分级构造密切相关（Haldane, 1983; Halperin, 1984）。

刻画 FQH 液体内部结构的首次尝试由 Girvin 和 MacDonald (1987) 提出，他们证明 Laughlin 态在一个非局域算符中具有非对角长程序。这一观察导致用 Ginzburg--Landau--Chern--Simons 有效理论来描述 FQH 液体（Read, 1989; Zhang et al., 1989）。这些进展产生了许多有趣而重要的结果，并加深了对 FQH 液体的理解。然而，这里我将尝试从更一般的观点出发描述 FQH 液体的内部结构。看来，FQH 液体的某些内部结构（特别是在所谓非阿贝尔 FQH 液体中的那些）无法用 Ginzburg--Landau--Chern--Simons 有效理论及相关的非对角长程序来描述。因此，我们需要发展更普遍的概念与表述，如拓扑序与拓扑场论，来描述 FQH 液体中的内部结构（见第 8 章）。在我看来，非对角长程序的概念并没有抓住 FQH 液体内部结构的本质。这里我们将不使用非对角长程序来发展有效理论（Blok and Wen, 1990a,b）。所得的有效理论只含纯 Chern--Simons 项。纯 Chern--Simons 形式更为紧凑，也更清楚地揭示分级 FQH 态的内部结构。

\subsection{Laughlin 态的有效理论}

\begin{keypoints}
\item FQH 液体的流体动力学与有效 Chern--Simons 理论。
\end{keypoints}

本节中我们只考虑单层自旋极化的 QH 系统。为构造分级态的有效理论，我们先设法得到 Laughlin 态的有效理论，然后在下一节用分级构造得到分级态的有效理论。

考虑磁场中的一个电子系统。为了涵盖更一般的情形，我们假设电子可以具有玻色统计或费米统计。系统的拉格朗日量具有如下形式（第一次量子化的版本）：
\begin{equation*}
\begin{aligned}
\mathcal{L}&=-eA^iJ^i+\text{动能/势能}\\
&=eA_iJ^i+\text{动能/势能}
\end{aligned}
\tag{7.3.1}
\end{equation*}
其中
\begin{equation*}
\pmb{J}(\pmb{x})=\sum_i\pmb{v}_i\delta(\pmb{x}-\pmb{x}_i),\qquad J^0(\pmb{x})=\sum_i\delta(\pmb{x}-\pmb{x}_i)
\tag{7.3.2}
\end{equation*}
分别是电子的流与密度，$(\pmb{x}_i,\pmb{v}_i)$ 是第 $i$ 个电子的位置与速度。动能/势能由 $\sum_i\frac{1}{2}mv_i^2+\sum_{i<j}V(\pmb{x}_i-\pmb{x}_j)$ 给出，其具体形式对我们并不重要。我们还假设了每个电子的电荷为 $-e$，光速为 $c=1$。

在流体动力学方法中，我们假设低能集体模式可以用密度与流 $J^\mu$ 描述，而低能有效理论具有形式 $\mathcal{L}(A_\mu,J^\nu)=eA_iJ^i+\mathcal{L}'(J^\mu)$。从 5.1.3 节与 5.3.3 节的讨论我们看到，有时一个态的低能激发可能太多，以至于无法用单一密度模式来描述。由于有能隙的 FQH 态的低能激发比超流体还要少，可以合理地假设单一密度模式足以描述 FQH 态的低能涨落。

在填充因子 $\nu=1/m$ 处（$m$ 对玻色型电子为偶整数，对费米型电子为奇整数），电子系统的基态由 Laughlin 波函数（Laughlin, 1983）给出：
\begin{equation*}
\Bigl[\prod(z_i-z_j)^m\Bigr]\mathrm{e}^{-\frac{1}{4l_B^2}\sum|z_i|^2}
\tag{7.3.3}
\end{equation*}
（这里我们已假设 $B<0$，于是 $(-e)B>0$。）为构造有效理论，我们注意到态 (7.3.3) 是不可压缩流体，其密度与磁场绑定，即 $J^0=(-e)\nu B/2\pi$。再结合有限的霍尔电导 $\sigma_{xy}=\frac{\nu e^2}{h}=\frac{\nu e^2}{2\pi\hbar}=\frac{\nu e^2}{2\pi}$，我们发现电子数流 $J^\mu$ 对电磁场的变化有如下响应：
\begin{equation*}
-e\delta J^\mu=-\sigma_{xy}\varepsilon^{\mu\nu\lambda}\partial_\nu\delta A_\lambda=-\frac{\nu e^2}{2\pi}\varepsilon^{\mu\nu\lambda}\partial_\nu\delta A_\lambda
\tag{7.3.4}
\end{equation*}
我们这样选择有效拉格朗日量 $\mathcal{L}(A_\mu,J^\nu)$，使它产生运动方程 (7.3.4)。为方便起见，引入一个 $U(1)$ 规范场 $a_\mu$ 来描述电子数流：
\begin{equation*}
J^\mu=\frac{1}{2\pi}\partial_\nu a_\lambda\varepsilon^{\mu\nu\lambda}
\end{equation*}
这样定义的流自动满足守恒定律。于是，产生式 (7.3.4) 的有效拉格朗日量取如下形式：
\begin{equation*}
\mathcal{L}=\Bigl[-m\frac{1}{4\pi}a_\mu\partial_\nu a_\lambda\varepsilon^{\mu\nu\lambda}+\frac{e}{2\pi}A_\mu\partial_\nu a_\lambda\varepsilon^{\mu\nu\lambda}\Bigr]
\tag{7.3.5}
\end{equation*}

式 (7.3.5) 只描述基态对外电磁场的线性响应。要对拓扑流体（例如 FQH 液体）有更完整的描述，我们需要把电子激发引入有效理论。特别是，我们要确保有效理论中包含一个携带与电子相同量子数的激发。

\subsection{有效理论中的电子与准粒子激发}

\begin{keypoints}
\item 单凭有效 Chern--Simons 理论并不是对 FQH 态的完整描述。为得到完整的有效理论，我们还需要指明电子算符。
\item 电子与准粒子/准空穴之间所要求的平凡互统计决定了准粒子与准空穴的量子数。
\end{keypoints}

让我们引入一个携带 $a_\mu$ 电荷 $l$ 的粒子。在有效理论 (7.3.5) 中，这样的粒子对应于如下源项：
\begin{equation*}
la_0\delta(\pmb{x}-\pmb{x}_0)
\tag{7.3.6}
\end{equation*}
该源项将产生一个电荷为
\begin{equation*}
Q=-le/m.
\tag{7.3.7}
\end{equation*}
的激发。这可以从运动方程 $\delta\mathcal{L}/\delta a_0=0$ 看出，由此方程得到
\begin{equation*}
J_0=\frac{1}{2\pi}\varepsilon_{ij}\partial_i a_j=-\frac{e}{2\pi m}B+\frac{l}{m}\delta(\pmb{x}-\pmb{x}_0)
\end{equation*}
右端第一项表明填充因子 $\nu\equiv 2\pi\frac{J_0}{-eB}$ 确实为 $\nu=1/m$，第二项则对应于与该激发相联系的电子密度的增加。

我们还看到，由源项 (7.3.6) 产生的激发带有 $l/m$ 个单位的 $a_\mu$ 通量。这样，如果有两个分别携带 $a_\mu$ 电荷 $l_1$ 与 $l_2$ 的激发，那么把一个激发绕另一个激发移动一周将诱导一个相位 $2\pi\times(a_\mu\text{-通量量子数})\times(a_\mu\text{ 电荷})$，\footnote{细心的读者可能注意到，这两个激发既带有 $a_\mu$ 电荷又带有 $a_\mu$ 通量。因此，把一个激发绕另一个激发移动一周所诱导的相位似乎应收到两项贡献：一项来自电荷绕通量运动，另一项来自通量绕电荷运动，如图 7.8 所描述的那样。这将给出相位 $2\pi\times\frac{l_1}{m}\times l_2+2\pi\times\frac{l_2}{m}\times l_1$，是式 (7.3.8) 中结果的两倍。然而，更仔细的计算表明，图 7.8 的图像并不适用于由 Chern--Simons 项诱导的电荷--通量束缚态。式 (7.3.8) 中的朴素结果恰好是正确的（见问题 7.3.1）。}即
\begin{equation*}
2\pi\times\frac{l_1}{m}\times l_2
\tag{7.3.8}
\end{equation*}
如果 $l_1=l_2\equiv l$，那么这两个激发是全同的。把它们互换将诱导式 (7.3.8) 中相位的一半，即
\begin{equation*}
\theta=\pi\frac{l^2}{m}
\tag{7.3.9}
\end{equation*}
这里的 $\theta$ 正是携带 $l$ 个单位 $a_\mu$ 电荷的激发的统计角。

我们的电子携带电荷 $-e$。从上面的讨论我们看到，电荷为 $-e$ 的激发对应于携带 $m$ 个单位 $a_\mu$ 电荷的粒子。这样的激发具有统计角 $\theta=\pi m$（见式 (7.3.9)）；因此，$m$ 为偶数时它是玻色子，$m$ 为奇数时它是费米子。于是，对玻色型电子取偶数 $m$ 的情形，以及对费米型电子取奇数 $m$ 的情形，我们可以把携带 $m$ 个单位 $a_\mu$ 电荷的激发认定为构成 FQH 液体的电子。具有电子的电荷与统计的激发的存在，支持了我们的 Chern--Simons 有效理论的正确性。

我们想强调，在有效理论中认定基本电子是非常重要的。正是这一认定，连同有效拉格朗日量，提供了对 FQH 液体拓扑性质的完整描述。我们将在下面看到，这一认定使我们能够确定准粒子激发的分数电荷与分数统计。

位置由复数 $\xi=x+\mathrm{i}y$ 描述的准空穴激发，是通过用 $\prod_i(\xi-z_i)$ 乘以基态波函数 (7.3.3) 产生的。注意到当一个电子绕准空穴一周时，波函数的相位改变 $2\pi$。这个 $2\pi$ 相位意味着电子与准空穴具有平凡的互统计。一般地，电子与一个允许的激发必须具有平凡的互统计，这样该激发的电子波函数才是单值的。

现在让我们试着通过插入一个携带 $l$ 个单位 $a_\mu$ 电荷的源项来产生激发。把一个电子绕这样的激发移动一周将诱导相位 $2\pi l$（见式 (7.3.8)）。电子波函数的单值性要求这样的相位是 $2\pi$ 的整数倍。因此 $a_\mu$ 电荷必须量子化为整数，也只有这些电荷对应的才是允许的激发。

从式 (7.3.7) 给出的激发电荷我们发现，$l=-1$ 对应于基本准空穴激发，而 $l=1$ 对应于基本准粒子激发。准粒子激发携带电荷 $-e/m$，准空穴携带电荷 $e/m$。从式 (7.3.9) 可见，二者都具有统计 $\theta=\pi/m$。我们看到，有效理论重现了 Laughlin 态中准粒子的著名结果（Arovas et al., 1984）。包含准粒子激发的完整有效理论为
\begin{equation*}
\begin{aligned}
\mathcal{L}&=\Bigl[-m\frac{1}{4\pi}a_\mu\partial_\nu a_\lambda\varepsilon^{\mu\nu\lambda}+\frac{e}{2\pi}A_\mu\partial_\nu a_\lambda\varepsilon^{\mu\nu\lambda}\Bigr]\\
&\quad+la_\mu j^\mu+\text{动能/势能}
\end{aligned}
\tag{7.3.10}
\end{equation*}
其中 $j^\mu$ 是准粒子的流，其形式由式 (7.3.2) 给出。式 (7.3.10) 连同对 $l$ 的量子化条件，是一个完整的低能有效理论，它捕捉了 $1/m$ Laughlin 态的拓扑性质。

\subsection*{问题 7.3.1}

证明式 (7.3.8)。（提示：可以从下面带两个激发的有效理论出发：
\begin{equation*}
\mathcal{L}=\Bigl[-\frac{1}{2}\frac{m}{2\pi}a_\mu\partial_\nu a_\lambda\varepsilon^{\mu\nu\lambda}+j^\mu a_\mu\Bigr]
\end{equation*}
其中 $j^\mu=j_1^\mu+j_2^\mu$，而
\begin{equation*}
\begin{aligned}
j_1^0(\pmb{x},t)&=l_1\delta(\pmb{x}-\pmb{x}_1(t)),&\qquad j_2^0(\pmb{x},t)&=l_2\delta(\pmb{x}-\pmb{x}_2(t)),\\
j_1^i(\pmb{x},t)&=l_1\dot{x}_1^i\delta(\pmb{x}-\pmb{x}_1(t)),&\qquad j_2^i(\pmb{x},t)&=l_2\dot{x}_2^i\delta(\pmb{x}-\pmb{x}_2(t)).
\end{aligned}
\end{equation*}
这里 $j^\mu$ 是两个激发的总流，$\pmb{x}_{1,2}(t)$ 是两个激发的位置。然后可以积掉 $a_\mu$，得到
\begin{equation*}
\int\mathrm{d}^3x\;\frac{1}{2}j^\mu\frac{1}{\frac{m}{2\pi}\partial_\lambda\epsilon^{\mu\lambda\nu}}j^\nu
\end{equation*}
交叉项可以写成
\begin{equation*}
\int\mathrm{d}^3x\;j_1^\mu\frac{1}{\frac{m}{2\pi}\partial_\lambda\epsilon^{\mu\lambda\nu}}j_2^\nu=\int\mathrm{d}^3x\;j_1^\mu f_\mu
\end{equation*}
其中 $f_\mu$ 满足
\begin{equation*}
\frac{m}{2\pi}\partial_\lambda\epsilon^{\mu\lambda\nu}f_\nu=j_2^\mu
\end{equation*}
假设 $\pmb{x}_2=\dot{\pmb{x}}_2=0$，即可求出 $f_i$。）

\subsection*{问题 7.3.2}

前面我们一直关注 Laughlin 态的拓扑性质。要对 Laughlin 态的动力学性质有所了解，我们需要在有效 Chern--Simons 理论中纳入如下的 Maxwell 项：
\begin{equation*}
\mathcal{L}=-m\frac{1}{4\pi}a_\mu\partial_\nu a_\lambda\varepsilon^{\mu\nu\lambda}+\frac{1}{2g_1}e^2-\frac{1}{2g_2}b^2
\tag{7.3.11}
\end{equation*}
其中 $e$ 与 $b$ 分别是 $a_\mu$ 的电场与磁场，且 $e_i=\dot{a}_i-\partial_i a^0$，$b=\partial_1a_2-\partial_2a_1$。求由 $e$ 与 $b$ 描述的集体涨落的运动方程。求集体激发的能隙。假设 $1/m$ Laughlin 态中的能隙为 $e^2/m^2 l_B$ 的量级，准空穴的尺寸为 $l_B$ 的量级。估计 $g_1$ 与 $g_2$ 的值。

\subsection{分级分数量子霍尔态的有效理论}

\begin{keypoints}
\item 不同的分级 FQH 态可以用一个整数 $K$ 矩阵和一个电荷矢量 $\pmb{q}$ 表征。
\item FQH 态的拓扑性质，如填充因子与准粒子量子数，可以容易地由 $K$ 矩阵和电荷矢量算出。
\item 不同的 $K$ 矩阵--电荷矢量对之间的等价关系。
\end{keypoints}

为得到分级 FQH 态的有效理论，让我们从费米型电子形成的 $1/m$ Laughlin 态出发。通过产生基本准粒子（记为 $l=1$）来增大填充因子。$l=1$ 的式 (7.3.10) 描述了存在这些准粒子时的 $1/m$ 态。现在出现两幅等价的图像。

\textbf{(a)}~在平均场理论方法中，我们可以把式 (7.3.10) 中的规范场 $a_\mu$ 视为固定背景，不允许 $a_\mu$ 对插入的源项 $j^\mu$ 作出响应。在这种情形下，准粒子气体的行为像"磁"场 $b=\partial_i a_j\varepsilon_{ij}$ 中的玻色子，这从式 (7.3.10) 的第二项可以看出。这些玻色子不携带任何电荷，因为准粒子数流 $j^\mu$ 不与电磁规范势 $A_\mu$ 直接耦合。当玻色子密度满足
\begin{equation*}
j^0=\frac{1}{p_2}\frac{b}{2\pi}
\end{equation*}
其中 $p_2$ 为偶数时，玻色子具有填充因子 $\frac{1}{p_2}$。玻色子的基态又可以用一个 Laughlin 态来描述。我们得到的最终电子态正是 Haldane (1983) 所构造的第二级分级 FQH 态。

\textbf{(b)}~如果我们让 $a_\mu$ 对 $j^\mu$ 的插入作出响应，那么准粒子将被 $a_\mu$ 通量缀饰。被缀饰的准粒子携带电荷 $-e/m$，统计为 $\theta=\pi/m$。当准粒子具有密度
\begin{equation*}
j^0=\frac{1}{(p_2-\frac{\theta}{\pi})}\frac{(-e)B}{2\pi m}
\end{equation*}
其中 $p_2$ 为偶数时，准粒子将具有填充因子 $\frac{1}{(p_2-\frac{\theta}{\pi})}$。这时，准粒子系统可以形成由波函数
\begin{equation*}
\prod_{i<j}(z_i-z_j)^{p_2-\frac{\theta}{\pi}}
\end{equation*}
描述的 Laughlin 态。用这种方法得到的最终电子态同样是一个第二级分级 FQH 态。这一构造由 Halperin (1984) 首先提出。(a) 与 (b) 两种构造给出同一个分级态，它们是等价的。

下面，我们将遵循 Haldane 的分级构造来推导分级 FQH 态的 Chern--Simons 有效理论（Blok and Wen, 1990a,b; Read, 1990; Fr\"{o}hlich and Kerler, 1991; Wen and Zee, 1992a）。注意到，在假设 (a) 之下，玻色子拉格朗日量（$l=1$ 时式 (7.3.10) 的第二项）正是把外电磁场 $eA_\mu$ 换成 $a_\mu$ 后的式 (7.3.1)。于是，我们可以重复从式 (7.3.3) 到式 (7.3.10) 的步骤，来构造玻色子 Laughlin 态的有效理论。引入一个新的 $U(1)$ 规范场 $\tilde{a}_\mu$ 来描述玻色子流，我们发现玻色子有效理论具有形式
\begin{equation*}
\mathcal{L}=-\frac{p_2}{4\pi}\tilde{a}_\mu\partial_\nu\tilde{a}_\lambda\varepsilon^{\mu\nu\lambda}+\frac{1}{2\pi}a_\mu\partial_\nu\tilde{a}_\lambda\varepsilon^{\mu\nu\lambda}
\tag{7.3.12}
\end{equation*}
在式 (7.3.12) 中，新的规范场 $\tilde{a}_\mu$ 描述玻色子的密度 $j^0$ 与流 $j^i$，它由
\begin{equation*}
j^\mu=\frac{1}{2\pi}\partial_\nu\tilde{a}_\lambda\varepsilon^{\mu\nu\lambda}
\end{equation*}
给出。这把 $a_\mu$ 与玻色子流 $a_\mu j^\mu$ 的耦合约化为 $a_\mu$ 与 $\tilde{a}_\mu$ 之间的一个 Chern--Simons 项（它成为式 (7.3.12) 中的第二项）。总的有效理论（包括原来的电子凝聚体）具有形式
\begin{equation*}
\begin{aligned}
\mathcal{L}&=\Bigl[-\frac{p_1}{4\pi}a_\mu\partial_\nu a_\lambda\varepsilon^{\mu\nu\lambda}+\frac{e}{2\pi}A_\mu\partial_\nu a_\lambda\varepsilon^{\mu\nu\lambda}\Bigr]\\
&\quad+\Bigl[-\frac{p_2}{4\pi}\tilde{a}_\mu\partial_\nu\tilde{a}_\lambda\varepsilon^{\mu\nu\lambda}+\frac{1}{2\pi}a_\mu\partial_\nu\tilde{a}_\lambda\varepsilon^{\mu\nu\lambda}\Bigr]
\end{aligned}
\tag{7.3.13}
\end{equation*}
其中 $p_1=m$ 是奇整数。式 (7.3.13) 就是第二级分级 FQH 态的有效理论。
